\subsection{2. Kết quả thực nghiệm và Đánh giá}

\subsubsection{2.1. Kết quả định lượng tổng thể}
\indent\indent Hiệu năng của mô hình đề xuất được so sánh với các kiến trúc cơ sở (Baselines) trên cùng tập kiểm thử 2,000 mẫu chuẩn của DTD-lightweight.

\begin{table}[H]
    \centering
    \caption{So sánh hiệu năng giữa các mô hình trên tập kiểm thử}
    \label{tab:model_comparison}
    \vspace{0.2cm}
    \renewcommand{\arraystretch}{1.3}
    \begin{tabular}{|m{5.5cm}|m{2.5cm}|m{2.5cm}|m{2.5cm}|}
        \hline
        \textbf{Mô hình} & \textbf{Validation Loss} & \textbf{Accuracy} & \textbf{F1-Score} \\
        \hline
        Baseline CNN & 0.6210 & 71.20\% & 0.5840 \\
        \hline
        ResNet-18 & 0.5124 & 74.30\% & 0.6120 \\
        \hline
        ResNet-50 & 0.4480 & 76.80\% & 0.6410 \\
        \hline
        \textbf{Ours (EfficientNet-B0 + SmolVLM2)} & \textbf{0.4011} & \textbf{78.50\%} & \textbf{0.6658} \\
        \hline
    \end{tabular}
\end{table}

Kết quả tại \textbf{Bảng \ref{tab:model_comparison}} cho thấy mô hình đề xuất đạt hiệu năng vượt trội trên mọi chỉ số: hàm mất mát kiểm định giảm xuống chỉ còn \textbf{0.4011}, độ chính xác tổng thể đạt \textbf{78.50\%} và điểm \textbf{F1-Score đạt 0.6658}, cao hơn đáng kể so với mạng ResNet-50 dù kích thước nhỏ hơn gần 5 lần.

\subsubsection{2.2. Đánh giá chi tiết trên từng lớp nhãn làm giả}
\indent\indent Bảng \ref{tab:per_class_results} thể hiện chi tiết điểm F1-Score trên từng đặc trưng pháp y:
\begin{table}[H]
    \centering
    \caption{Hiệu năng phân loại chi tiết theo từng lớp nhãn}
    \label{tab:per_class_results}
    \vspace{0.2cm}
    \renewcommand{\arraystretch}{1.3}
    \begin{tabular}{|m{1cm}|m{5.5cm}|m{2.5cm}|m{2.5cm}|m{2cm}|}
        \hline
        \textbf{STT} & \textbf{Nhãn làm giả} & \textbf{Precision} & \textbf{Recall} & \textbf{F1-Score} \\
        \hline
        1 & Edge artifacts & 0.768 & 0.718 & \textbf{0.742} \\
        \hline
        2 & Unnatural texture & 0.702 & 0.669 & \textbf{0.685} \\
        \hline
        3 & Inconsistent font & 0.735 & 0.687 & \textbf{0.710} \\
        \hline
        4 & Inconsistent alignment & 0.654 & 0.623 & \textbf{0.638} \\
        \hline
        5 & Lack of integration & 0.641 & 0.608 & \textbf{0.624} \\
        \hline
        6 & Text incoherence & 0.615 & 0.578 & \textbf{0.596} \\
        \hline
        \multicolumn{2}{|c|}{\textbf{Trung bình (Macro Average)}} & \textbf{0.686} & \textbf{0.647} & \textbf{0.6658} \\
        \hline
    \end{tabular}
\end{table}

Các đặc trưng hình học và biên viền như \textit{Edge artifacts} (F1: 0.742) và \textit{Inconsistent font} (F1: 0.710) đạt độ chính xác cao nhất do mạng EfficientNet-B0 bắt rất nhạy các biến đổi đột ngột về độ dốc gradient và hoa văn font chữ. Nhãn \textit{Text incoherence} có điểm F1 thấp hơn (0.596) do đòi hỏi khả năng hiểu sâu về ngữ cảnh ngôn ngữ cấp tài liệu.

\subsubsection{2.3. Đánh giá định tính: Phân tích các trường hợp suy luận thực tế (Reasoning Case Studies)}
\indent\indent Dưới đây là ví dụ minh họa kết quả trích xuất suy luận từ mô hình SmolVLM2 trên một mẫu hóa đơn bị chỉnh sửa con số:
\begin{lstlisting}[language=TeX, frame=single, caption={Ví dụ kết quả suy luận trích xuất từ SmolVLM2}]
<think>
1. Inspect the highlighted region at line "Total Amount":
2. Digit '5' exhibits a distinct rectangular boundary discontinuity (Edge artifacts).
3. The stroke width and ink density of '5' differ from adjacent zeroes (Inconsistent font).
4. Local JPEG compression noise around '5' is unnaturally smooth compared to paper background (Unnatural texture).
Conclusion: The text area is tampered using text splicing manipulation.
</think>
\end{lstlisting}
Chuỗi tư duy trên chứng minh mô hình không chỉ dự đoán đúng nhãn mà còn chỉ ra chính xác các cơ sở vật lý vi mô, mang lại giá trị giải trình pháp lý cao.
